% ===========================================================================
\chapter{Setting up}\label{ch:setup}

Make a folder for the project and go into it. Every file in this tutorial
goes in this one folder.

\begin{shell}
mkdir island
cd island
\end{shell}

Copy the map, \code{kestrel_island.map}, from the \code{examples/digital_map}
folder of the Nex distribution into it. It is a plain text file; we look
inside it in Chapter~\ref{ch:mapfile}.

\section{A first file: helpers}

The first file holds a few small helpers that the rest of the project uses
for printing.

First, a word about printing itself. Nex's \code{print} shows a string
with its quotes, as \code{"hello"}, which is handy when debugging but not
in a program's output. Every Nex program has a \code{Console} called
\code{console}, ready to use without creating it, and
\code{console.print_line("hello")} prints the text as it is, on a line of
its own. We use it for all our output.

Now numbers. Nex prints a \code{Real} with all its digits, like
\code{1879.3774630534463}, which is far more than anyone wants to read. The
built-in \code{format} writes a number with a fixed number of decimal
places, as text: \code{2.5.format(3)} is \code{"2.500"}, and
\code{(-3.14159).format(2)} is \code{"-3.14"}. It rounds to the nearest,
and a number exactly halfway rounds away from zero.

A navigator doesn't say ``in 347.2 metres''. It rounds, and the less
precise the better, as long as it is precise enough: ``in 350 m'',
``in 1.2 km''. Under a kilometre we round to the nearest 10 metres, which
is what \code{(metres / 10.0).round * 10} does; from a kilometre up we
give kilometres with one decimal. Times are given in whole minutes.

\begin{newfile}{helpers.nex}
-- helpers.nex: printing distances and times neatly, and finding the
-- user's files.

-- A distance the way a navigator says it: "40 m", "350 m", "1.2 km".
-- Under a kilometre it is rounded to 10 m, which is as precise as a
-- driver needs.
function distance_text(metres: Real): String
require
  non_negative: metres >= 0.0
do
  if metres < 995.0 then
    result := ((metres / 10.0).round * 10).to_string + " m"
  else
    result := (metres / 1000.0).format(1) + " km"
  end
end

-- A duration in whole minutes, "under a minute" or "1 h 05 min".
function duration_text(seconds: Real): String
require
  non_negative: seconds >= 0.0
do
  let minutes: Integer := (seconds / 60.0).round
  if minutes < 1 then
    result := "under a minute"
  elseif minutes < 60 then
    result := minutes.to_string + " min"
  else
    result := (minutes / 60).to_string + " h "
      + (minutes % 60).to_string.pad_start("0", 2) + " min"
  end
end
\end{newfile}

Why does \code{distance_text} switch to kilometres at 995 metres rather
than 1,000? Because 997 metres rounds to ``1000 m'', which looks odd; from
995 up we say ``1.0 km'' instead.

\code{minutes / 60} divides two \code{Integer}s, so it gives an
\code{Integer}, dropping the remainder: 65 minutes is 1 hour, and
\texttt{65~\%~60}, the remainder, is the 5 minutes left over. The minutes
always get two digits, \code{05} rather than \code{5}: the built-in
\code{pad_start} does that, as \code{s.pad_start("0", 2)} adds \code{"0"}s
on the left of \code{s} until it is 2 characters long.

The last helper finds a file the user names on the command line. The
\code{nex} launcher runs programs from its own installation folder, so a
plain name like \code{kestrel_island.map} would be looked for in the wrong
place. Every program also has \code{program}, the running
\code{Program}, which knows the folder you ran \code{nex} from:
\code{program.working_directory}. We put that in front of any name that
isn't already a full path:

\begin{append}{helpers.nex}
-- The full path of a file the user named. The `nex` launcher runs
-- programs from its own installation folder, so a relative name is
-- resolved against the folder the program was started from.
function user_path(raw: String): String do
  if raw.starts_with("/") then
    result := raw
  else
    result := program.working_directory + "/" + raw
  end
end
\end{append}

\section{Running a try-it script}

Let's check that it works. Make this file and run it:

\begin{tryit}{try_helpers.nex}
intern Helpers

console.print_line(2.5.format(3))
console.print_line((-3.14159).format(2))
console.print_line(distance_text(347.2) + ", " + distance_text(997.0) + ", "
  + distance_text(1234.0))
console.print_line(duration_text(25.0) + ", " + duration_text(150.0) + ", "
  + duration_text(3900.0))
\end{tryit}
\tryoutput{try_helpers}

\code{intern Helpers} finds \code{helpers.nex} in the same folder and loads
everything in it. The file does not need a class called \code{Helpers}:
ours holds only plain functions.

Try-it scripts are not part of the project. Keep them around to experiment
with, or delete them when you are done.
